% Standalone section -- paste into the weekly report, or \input it.
% Week 3 working note: closes the open question left at the end of
% tortuosity_and_representation.tex ("A quantised-synthetic-curve test would settle the noise
% floor of each measure"). Companion to tortuosity_and_disease.tex, which asks whether the
% reproducible metrics are the discriminating ones; this one asks the prior question, how much of
% the derivative-based measures is signal at all, at the smoothing scale the rest of the panel
% already uses by default.
%
% PROVENANCE. scripts/tortuosity_noise_floor.py, run 2026-09-16, 800 usable cases plus case 1 for
% the figure. Figure from scripts/make_noise_floor_figure.py. Both are my own analysis, not a
% literature result, in the same spirit as the T_kappa noise-bias derivation in
% tortuosity_and_representation.tex Sec.\ 4.

\section*{Tortuosity}

\textbf{Tortuosity} isn't one thing. The word bundles together several distinct geometric
properties: how much longer a path is than the straight-line distance ($L/D$), how sharply it
bends locally (curvature), how much it corkscrews out of a flat plane (torsion, non-planarity),
how often it oscillates (inflections, SOAM). \\
The measures that matter most, curvature and torsion, are also the least trustworthy, because
they need derivatives. A centerline is a list of points with some position error in it: jitter
from skeletonisation, rounding to the voxel grid. Curvature takes two derivatives of that noisy
signal, torsion takes three, and each derivative makes the noise worse. \\

I started by investigating: at the current smoothing level, how much of the reported curvature
is real bending, and how much is just noise?


\subsection*{Lower bound for the noise}

A centerline point comes from a segmentation mask, which is a grid of voxels: each voxel is
either ``vessel'' or not. Wherever the true, continuous vessel boundary actually falls inside a
voxel, the mask can only place it at the voxel's resolution, not any finer. This rounding is
\emph{quantisation}, and it has a known, standard result: if a true position is rounded onto a
grid of spacing $h$, and there is no reason for that position to line up with the grid, i.e.\ a
patient's coronary anatomy has no relationship to the CT scanner's reconstruction grid, the
rounding error has variance exactly $h^2/12$, i.e.\ a standard deviation of
\begin{equation}
  \varepsilon = \frac{h}{\sqrt{12}}.
  \label{eq:week3-quant-floor}
\end{equation}
Any further source of error, segmentation mistakes, skeletonisation artefacts, only adds more
noise on top of this; it cannot remove it. So the total position noise is always
\textbf{at least} $h/\sqrt{12}$.

Reading the actual NIfTI headers of 60 ImageCAS-X volumes gives:
\begin{itemize}
  \item \textbf{in-plane} ($x$, $y$): 0.309--0.430~mm, varies scan to scan (median 0.359~mm)
  \item \textbf{through-plane} ($z$): exactly 0.5~mm, on every one 
\end{itemize}
I chose to take the $z$ spacing to be sure I don't understate the true floor on any scan:
\[
  \varepsilon \;\geq\; \frac{0.5~\text{mm}}{\sqrt{12}} \;\approx\; 0.144~\text{mm}.
\]

This noise-floor is the best case for the data, since we must expect some humen error when tracing the arteries, which is not added her.

\subsection{The derivative problem}


\textbf{Why derivatives make noise worse, in one example.} Take three points on a perfectly
straight line ($h = 0.5$~mm apart), each off by the position noise derived in
position\_noise\_lower\_bound.tex ($\varepsilon = 0.144$~mm), alternating sign:
$y_{i-1}, y_i, y_{i+1} = +\varepsilon, -\varepsilon, +\varepsilon$. The tangent,
$(y_{i+1}-y_{i-1})/2h$, comes out around $\varepsilon/h$: noise in, noise out, roughly the same
size. The curvature, $(y_{i+1} - 2y_i + y_{i-1})/h^2 = 4\varepsilon/h^2 \approx 2.3$~mm$^{-1}$, is
ten times larger than a real coronary bend ($0.05$--$0.3$~mm$^{-1}$) from the same $0.144$~mm
wiggle. One extra derivative, and a division by $h^2$ instead of $h$, turns invisible position
noise into a curvature reading bigger than any real vessel produces.


So computing curvature or torsion point to point is out of the question. Noise is roughly
zero-mean and independent point to point, so averaging $N$ points shrinks the noise variance by
about $1/N$, while a real bend, which persists over many consecutive points rather than flipping
sign from one to the next, does not average away the same way. Averaging is the way out; the only
questions left are how much of it, and measured against what.

\subsection*{The floor, measured}

Not the three-point sketch above, but the same idea done properly: 500 straight lines, each
150~mm long, every point given independent noise of size $\varepsilon = 0.144$~mm. Each line is
smoothed and differentiated exactly as a real vessel is, at a few different smoothing scales
$\sigma$. Since every one of these lines is truly straight, every curvature or torsion value that
comes out the other end is pure noise; the 95th percentile of that pooled result, one per
$\sigma$, is the floor.

\begin{table}[htbp]
  \centering
  \small
  \begin{tabular}{@{}rrr@{}}
    \toprule
    Smoothing $\sigma$ (mm) & curvature floor (mm$^{-1}$) & torsion floor (mm$^{-1}$) \\
    \midrule
    0.0 & 3.218 & 11.790 \\
    0.5 & 0.760 & 8.441 \\
    1.0 & 0.169 & 5.225 \\
    2.0 & 0.031 & 2.754 \\
    \bottomrule
  \end{tabular}
  \caption{The noise floor at each smoothing scale: the curvature/torsion value a perfectly
    straight line still reports, from 500 noisy synthetic lines
    (\texttt{scripts/tortuosity\_noise\_floor.py}). More smoothing lowers the floor a lot, but
    also erases real bends, so it is not free (\Cref{tab:week3-sweep} in
    tortuosity\_and\_representation.tex).}
  \label{tab:week3-noise-floor}
\end{table}

\subsection*{What survives at $\sigma = 1.0$mm}

Take a real vessel's own curvature at every point along its length, at the codebase's default
smoothing, and check how much of it clears the matching floor above. Across the 800 usable cases,
it is low everywhere: the LAD clears its own floor for a median of 14.5\% of its length, the LCx
for 12.6\%, the RCA for only 4.3\%, and torsion for about 1.2\% on every vessel, so torsion is
essentially never trustworthy at this smoothing scale.

\begin{table}[htbp]
  \centering
  \small
  \begin{tabular}{@{}lrr@{}}
    \toprule
    Vessel & curvature above floor (median \% of length) & torsion above floor (median \%) \\
    \midrule
    LM  & 0.0\%  & 0.0\% \\
    LAD & 14.5\% & 1.2\% \\
    LCx & 12.6\% & 1.2\% \\
    RCA & 4.3\%  & 1.2\% \\
    \bottomrule
  \end{tabular}
  \caption{Fraction of each vessel's length whose curvature (torsion) is above the noise floor, at
    $\sigma = 1.0$~mm, across 800 usable cases.}
  \label{tab:week3-clearance}
\end{table}

\Cref{fig:week3-noise-floor} shows one example, case 1's RCA: the left panel is the floor itself
(curvature measured on a straight line, so all of it is noise); the right panel is the real
vessel, same floor line drawn on top. Only the shaded part above the line counts as real bending.

\begin{figure}[htbp]
  \centering
  \includegraphics[width=\linewidth]{content/background/figures/tortuosity_noise_floor_1.pdf}
  \caption{Left: curvature measured on 500 noisy straight lines at $\sigma=1.0$~mm; every value
    here is noise. Right: case 1's real RCA curvature profile, same floor line; only the shaded
    11.5\% of its length sits above it (\texttt{scripts/make\_noise\_floor\_figure.py}).}
  \label{fig:week3-noise-floor}
\end{figure}

\subsection*{Consistency with the existing pipeline}

\texttt{curvature\_profile}, the point-by-point function introduced for this test, is not used
elsewhere in the codebase. \texttt{topology.tortuosity.measure}, the function through which every
real vessel is processed, does not return that raw series: it reports a mean, an RMS, and a
windowed maximum over a sliding 20~mm stretch (\texttt{curvature\_rms\_window\_max}), the
averaging argued for above, implemented before this test was written. The point-by-point profile
serves here to calibrate the noise floor against which those existing averages should be checked,
at whatever $\sigma$ they were computed at.